% Neoficiální startovací šablona závěrečné práce FP VUT.
% Údaje konkrétní práce vyplňte v settings.tex.
\documentclass[zadani]{fpthesis}

% Písmo a odkazy zachované z výchozí šablony FIT.
\renewcommand{\rmdefault}{lmr}
\renewcommand{\sfdefault}{qhv}
\renewcommand{\ttdefault}{lmtt}
\csdoublequotesoff
\usepackage[unicode,colorlinks,plainpages=false]{hyperref}
\definecolor{template-ref}{RGB}{223,52,30}
\definecolor{template-cite}{HTML}{2F8F00}
\definecolor{template-url}{HTML}{092EAB}
% Titulní list převádí obsah na verzálky a očekává i název barvy RED.
\colorlet{RED}{red}
\hypersetup{linkcolor=template-ref,citecolor=template-cite,urlcolor=template-url}

\graphicspath{{figures/}}
\newcommand{\placeholder}[1]{\textcolor{RED}{\textbf{DOPLŇTE: #1}}}

% Tento soubor je hlavní místo pro údaje konkrétní práce.
% Typ změňte na project={BP} pro bakalářskou práci nebo project={DP} pro diplomovou.
\projectinfo{
  project={DP},
  faculty={FP},
  department.cs={\textcolor{RED}{Doplňte oficiální český název ústavu}},
  department.en={\textcolor{RED}{Add the official English department name}},
  department={DEF},
  title.cs={\textcolor{RED}{Doplňte český název práce}},
  title.en={\textcolor{RED}{Add the English thesis title}},
  title.length={14.5cm},
  author.name={\textcolor{RED}{Jméno}},
  author.surname={\textcolor{RED}{Příjmení}},
  author.title.p={},
  author.title.a={},
  supervisor.name={\textcolor{RED}{Jméno}},
  supervisor.surname={\textcolor{RED}{Příjmení vedoucího}},
  supervisor.title.p={},
  supervisor.title.a={},
  year={\textcolor{RED}{Doplňte rok}},
  date=\today,
  abstract.cs={\placeholder{český abstrakt}},
  abstract.en={\placeholder{anglický abstrakt}},
  keywords.cs={\placeholder{česká klíčová slova}},
  keywords.en={\placeholder{anglická klíčová slova}},
  declaration={\placeholder{závazné znění prohlášení podle aktuálních pravidel školy}},
  % acknowledgment={Text poděkování}, % nepovinné
  % extendedabstract={Text rozšířeného abstraktu}, % pouze je-li vyžadován
  % extabstract.odd={true}, % nepovinně začít rozšířený abstrakt na liché straně
}


\begin{document}
\hypersetup{pageanchor=false}
\maketitle
\tableofcontents
\clearpage
\hypersetup{pageanchor=true}

\chapter{Úvod}

\placeholder{Nahraďte tuto kapitolu vlastním úvodem a upravte členění práce
podle skutečného zadání a požadavků fakulty.}

\chapter{Technické ukázky}
\label{chap:examples}

Tato kapitola předvádí základní prvky sazby. Při psaní ji nahraďte vlastním
textem. Kniha může být citována samostatně \cite{latex-companion} a později
znovu \cite{latex-companion}. Webová dokumentace je uvedena v~\cite{latex-docs}.
Více zdrojů lze odkázat jedním příkazem \cite{knuth1984,grobid2009}.

\section{Odstavec a seznam}

Odstavce oddělujte prázdným řádkem. Krátký výčet lze zapsat například takto:
\begin{itemize}
  \item první položka,
  \item druhá položka,
  \item poslední položka.
\end{itemize}

\section{Rovnice a odkaz}

Rovnice~\eqref{eq:example} má vlastní štítek, na který lze odkazovat z~textu.
\begin{equation}
  a^2 + b^2 = c^2
  \label{eq:example}
\end{equation}

\section{Tabulka}

Tabulka~\ref{tab:example} ukazuje popisek a odkaz.
\begin{table}[ht]
  \centering
  \caption{Ukázková struktura tabulky}
  \label{tab:example}
  \begin{tabular}{ll}
    \toprule
    Sloupec A & Sloupec B \\
    \midrule
    položka 1 & hodnota 1 \\
    položka 2 & hodnota 2 \\
    \bottomrule
  \end{tabular}
\end{table}

\section{Obrázek}

Obrázek~\ref{fig:example} používá pouze rámeček, aby šablona nepotřebovala
další ukázkový binární soubor.
\begin{figure}[ht]
  \centering
  \fbox{\rule{0pt}{35mm}\rule{0.65\textwidth}{0pt}}
  \caption{Místo pro ukázkový obrázek}
  \label{fig:example}
\end{figure}

% Vlastní soubor vložte do figures/ a rámeček výše nahraďte například:
% \includegraphics[width=0.8\textwidth]{nazev-obrazku.pdf}

\chapter{Závěr}

\placeholder{Nahraďte tuto kapitolu vlastním závěrem odpovídajícím skutečně
provedené práci.}


\bibliographystyle{czplain}
\bibliography{references}

% Volitelné seznamy patří až za použitou literaturu.
% \listoffigures
% \listoftables
% \chapter*{Seznam použitých zkratek a symbolů}

\appendix
\chapter{Ukázková příloha}
\label{app:example}

Tato označená příloha ověřuje pořadí částí dokumentu. Před odevzdáním ji
nahraďte skutečnou přílohou, nebo ji spolu s~příkazem v~\texttt{main.tex}
odstraňte.


\end{document}
